\section{A descriptor that also carries its own class: the 639-bit packet}
\label{sec:e639}

A further representation change removes six term nodes from the closure
axiom.  Combined with the preceding truth-equation deletion and the shorter
ground query, it gives a complete \textbf{639-bit equational consistency
packet}: seven equations, nine operation symbols, four variable slots, and
144 equation-term nodes.  The query is \(Q(d)=d\).  The signature and every
equation and endpoint are included in these 639 logical bits; the physical
file occupies 80 bytes with one zero padding bit.

The change is to use the descriptor itself as the class code.  The original
692-bit packet represents its schema matrix by
\[
 R_b^{\mathrm{old}}(x,y)
   =p_\varepsilon(P(y,P(x,b)),V(b,d)).
\]
The new packet uses
\begin{equation}
                  R_b(x,y)=p_\varepsilon(P(y,P(x,b)),b).
\label{eq:e639-R}
\end{equation}
There are exactly three occurrences of \(V(b,d)\) in the old closure row.
Replacing each by \(b\) saves two term nodes, hence 24 bits.  The head
equation itself is retained.  A protected parameter tail proves that this
change still represents every matrix with every finite tuple of parameters.
The comparison below is a relative model theorem, not a claim that the old
and new closure equations have identical models.

\subsection{The literal seven-equation presentation}

The ordered signature remains
\[
 d^0,N^2,V^2,P^2,A^2,Q^1,W^2,p_E^2,p_\varepsilon^2.
\]
Write \(\varepsilon=p_\varepsilon\).  The first six equations are
\begin{align*}
 N(N(x,N(N(y,x),x)),N(y,N(z,x)))&=y,\\
 V(P(x,s),d)&=x,\\
 V(P(x,s),Q(j))&=V(s,j),\\
 \varepsilon(s,P(a,b))&=N(\varepsilon(s,a),\varepsilon(s,b)),\\
 \varepsilon(s,A(i,j))&=p_E(V(s,i),V(s,j)),\\
 \varepsilon(s,Q(a))
   &=N\bigl(\varepsilon(P(W(s,a),s),a),
            \varepsilon(P(z,s),a)\bigr).
\end{align*}
The last equation uses the following abbreviations for actual terms:
\begin{align*}
 H_b&=Q(b), & n_z&=P(b,z), & m_z&=A(b,z),\\
 r&=R_b(H_b,n_z), & e&=p_E(m_z,n_z), & f&=p_E(m_z,z),\\
 a&=p_E(W(x,d),x), & b_0&=p_E(x,H_b), & c&=R_b(x,y),\\
 u&=p_E(y,H_b), & v&=R_b(n_y,y).&&
\end{align*}
Define the terms
\[
 K=N(e,N(d,f)),\qquad
 L^-=N(c,N(a,N(d,b_0))),\qquad
 H=N\bigl(N(N(d,u),L^-),N(u,v)\bigr).
\]
The seventh equation is
\begin{equation}
                 d=N\bigl(N(r,N(d,K)),N(d,H)\bigr).
\label{eq:e639-closure}
\end{equation}
The file \texttt{E639.equations} expands every abbreviation.  No arity,
definition, or omitted helper is supplied outside the packet.

\subsection{The protected parameter tail}

Sh1 gives a Boolean reduct independently of the value of \(d\).  By the
class-NAND equation, the ground terms
\[
                    T=P(d,P(d,d)),\qquad F=P(T,T)
\]
are respectively a uniformly true class and a uniformly false class:
\[
 \varepsilon(s,T)=1,\qquad \varepsilon(s,F)=0
 \quad\hbox{for every environment }s.
\]
These are statements about predicate values.  They do not identify the
native object \(T\) with the Boolean top or \(F\) with the Boolean bottom.

Fix arbitrary native objects \(p_1,\ldots,p_k\), and set
\[
 L=P(p_1,P(p_2,\ldots P(p_k,d)\ldots)),\qquad
 B=P(c,P(F,L)).
\]
For \(k=0\), use \(L=d\).  The tail \(P(F,L)\) is a true class regardless
of the meanings of the parameters, because
\(\varepsilon(s,P(F,L))=N(0,\varepsilon(s,L))=1\).
Consequently
\begin{equation}
                       \varepsilon(s,B)=\neg\varepsilon(s,c).
\label{eq:e639-guard}
\end{equation}
This is the central coding device.  Merely putting parameters after the
code would not have this property: arbitrary parameter objects can have
arbitrary class meanings.  The fixed false-class head makes the parameter
tail semantically harmless without restricting any parameter.

Let \(\nu_0=d\) and \(\nu_{i+1}=Q(\nu_i)\).  The head/tail equations
show by induction that \(V(s,\nu_i)\) reads coordinate \(i\) of every
sufficiently long finite \(P\)-stack.  In
\(S=P(y,P(x,B))\), these coordinates are
\[
               (y,x,c,F,p_1,\ldots,p_k).
\]
Thus output has coordinate 0, input coordinate 1, the two internal
descriptor fields occupy coordinates 2 and 3, and the parameters start
at coordinate 4.  No assumption about the default tail is needed for
these finitely many coordinates.

After the usual Boolean homomorphism to two truth values, the retained
binder means existential negation: \(Q(a)\) represents
\(\exists z\,\neg a\).  Ordinary existential quantification is
\(Q(P(a,a))\).  Membership atoms are represented by
\(A(\nu_i,\nu_j)\), NAND by \(P\), and bound coordinates are shifted
on prepending the bound variable.  The standard structural induction
therefore represents every finite equality-free pure-membership matrix.
Given \(R(x,y;p_1,\ldots,p_k)\), first construct a class code \(c\) for
\(\neg R\) with the coordinate convention just stated, and only then
form \(B\).  By \eqref{eq:e639-guard}, the new matrix \(R_B\) is exactly
\(R\).  The code is built before its descriptor; there is no circular
definition and no restriction to parameter-free instances.

\subsection{Truth recovery and the ground query remain valid}

The omitted truth equation still follows.  In
\eqref{eq:e639-closure}, set \(b=T\).  Then
\(r=\varepsilon(P(P(T,z),P(Q(T),T)),T)=1\).
The same NAND identity used for E667 applies:
\[
 N\bigl(N(N(t(w),N(d,K)),N(d,H)),N(d,d)\bigr)=d,
 \qquad t(w)=N(w,N(w,w)).
\]
It gives \(N(d,N(d,d))=d\).  The native head equation is no longer even
needed for this particular deduction.  Hence \(d\) is the Boolean top
in every model of the seven retained equations.

The query \(Q(d)=d\) holds only in trivial models, by the head/tail chain
\[
 x=V(P(x,P(y,d)),d)=V(P(x,P(y,d)),Q(d))=V(P(y,d),d)=y.
\]
Conversely \(N(d,d)=d\) collapses every Sh1 Boolean reduct and implies
\(Q(d)=d\).  Thus the shortened query and the original Boolean-collapse
query have the same equational refutability for this new presentation too.

\subsection{From a separating model to raw EPSC}

Recall the full equality-free raw EPSC schema.  For every finite
pure-membership relation formula \(R(x,y;\bar p)\), if
\[
 \forall x\,\exists z\,\forall y\,
       \bigl(R(x,y;\bar p)\longrightarrow
          \forall m\,(E(m,y)\to E(m,z))\bigr),
\]
then some \(h\) satisfies both
\begin{align*}
 &\forall x,y\,
    \bigl((\mathsf{Empty}(x)\lor E(x,h))\land R(x,y;\bar p)
       \longrightarrow E(y,h)\bigr),\\
 &\forall y\,(E(y,h)\longrightarrow\exists x\,R(x,y;\bar p)),
 \qquad \mathsf{Empty}(x)\equiv\forall m\,\neg E(m,x).
\end{align*}

Take a model separating \(Q(d)\) and \(d\).  Its Boolean reduct is
nontrivial and its \(d\) is top.  Choose an ultrafilter of that Boolean
algebra and apply its homomorphism to \(\{0,1\}\) only to the values of
\(p_E\) and \(\varepsilon\).  Retain the entire native carrier and all
native operations.  In particular native stack equality remains literal
equality.  The selected complemented-witness equation gives
\[
 \varepsilon(s,Q(a))=1
       \quad\Longleftrightarrow\quad
       \exists z\,\varepsilon(P(z,s),a)=0.
\]
For completeness, set \(w=\varepsilon(P(W(s,a),s),a)\).  The instance
\(z=W(s,a)\) gives the binder value \(\neg w\), and its other instances
give \(w\leq\varepsilon(P(z,s),a)\).  In two truth values this is exactly
the displayed equivalence, including an attained witness.  No preservation
of arbitrary infinite joins by the homomorphism is assumed.

Let \(E(x,y)\) mean that the image of \(p_E(x,y)\) is true.  Every matrix
and parameter tuple is represented by a protected descriptor \(B\).
The closure equation becomes
\[
 \mathsf{Bad}(B,z)\lor\mathsf{Hull}(B,x,y),
\]
where
\begin{align*}
 \mathsf{Bad}&=r\land e\land\neg f,\\
 \mathsf{Hull}&=(a\lor\neg c\lor u)
                \land(\neg b_0\lor\neg c\lor u)
                \land(\neg u\lor v).
\end{align*}
If the represented relation is locally bounded, apply the bound at
\(H_B=Q(B)\).  Some \(z\) makes \(\mathsf{Bad}(B,z)\) false: otherwise
\(R_B(H_B,n_z)\), \(E(m_z,n_z)\), and \(\neg E(m_z,z)\) contradict
that bound.  Hence all hull clauses hold.  If \(x\) is empty then
\(a=E(W(x,d),x)\) is false, and the first clause gives closure at empty
inputs.  The second gives closure at members of the hull.  The third gives
unrestricted support, with witness \(P(B,y)\).  Thus every instance of
the full raw EPSC schema holds on the same carrier.

\subsection{The reverse model construction and its recursion order}

For the reverse implication, let \((M,E_M)\) be a countable nonempty
raw EPSC structure.  Use the same free carrier as the inspected E692
construction: finite trees over constructors \(P,A,Q\), with a leaf
\(d\) and separate leaves \(e_m\) for all \(m\in M\).  All these
native constructors are literal disjoint tree constructors.  The extra
leaves are carrier elements, not extra signature constants.

Define \(\operatorname{idx}(i)\) as the number of outer \(Q\)-layers
of the tree \(i\).  Read a \(P\)-tree as a cons cell and a non-\(P\)
tail as the constant-\(d\) stream; let \(V(s,i)\) read coordinate
\(\operatorname{idx}(i)\).  Both native stack equations hold for all
arguments, canonical or otherwise.  Decode a class structurally: leaves
are false, \(P\) is NAND, \(A(i,j)\) is membership between the selected
coordinates, and \(Q(a)\) means existential negation of the decoded
body after prepending its bound value.

Give leaves weight 1 and constructors weights
\[
 w(Q(b))=w(b)+1,\quad
 w(P(b,z))=w(b)+w(z)+2,\quad
 w(A(b,z))=w(b)+w(z)+3.
\]
For each descriptor \(b\), decode \emph{the class tree \(b\) itself}
on the stream \(P(y,P(x,b))\).  This gives a finite pure-membership
formula \(\widehat R_b(x,y;\bar p)\).  Each fixed parameter in that
formula is a proper tail-head subtree of \(b\) or the already assigned
default \(d\): the first two free stream coordinates are \(y,x\),
and subsequent coordinates are read from the native tail \(b\).
If that tail is not a \(P\)-tree, they are all default \(d\).
Bound variables inserted by class quantifiers do not change this
classification.  The class-decoding operation inspects finite syntax and
selector indices; it never queries an unassigned raw value.

This is the revised descriptor-descent argument.  In particular, evaluating
the class \(b\) does not amount to a truth predicate for the expanded
language, and it does not introduce a raw function term such as
\(\mu(P(x,b))\) into the matrix.  The matrix remains in the pure
membership language, with finitely many fixed parameters of smaller
weight.  This is what allows the raw EPSC hypothesis to be used.

Choose \(m_0\in M\).  Set \(\mu(d)=m_0\), \(\mu(e_m)=m\), and
define the remaining raw values recursively by weight.  Once the
parameters of \(\widehat R_b\) have their raw values, write its
resulting relation on \(M\) as \(R_b^M\).  If this relation is locally
bounded, choose an EPSC hull \(h_b\).  Otherwise choose \(h_b\) with
\[
 \forall z\,\exists n,m\,
      \bigl(R_b^M(h_b,n)\land E_M(m,n)\land\neg E_M(m,z)\bigr),
\]
the classical negation of local boundedness.  Put \(\mu(Q(b))=h_b\).
In the latter case choose \(\mu(P(b,z))\) as the displayed \(n\)
for the bound \(\mu(z)\), and choose \(\mu(A(b,z))\) as its displayed
separating member \(m\).  In the hull case, if
\(E_M(\mu(z),h_b)\), use support to choose \(\mu(P(b,z))=n\) with
\(R_b^M(n,\mu(z))\); otherwise use \(m_0\).  In that case set
\(\mu(A(b,z))=m_0\).

Every choice is available when needed.  Specifically,
\[
                   w(Q(b))<w(P(b,z))<w(A(b,z)),
\]
and every descriptor parameter has smaller weight than \(Q(b)\).
A fixed enumeration of the countable domain makes the choices consistent.
The leaves \(e_m\) make \(\mu\) surjective.

Pull membership back by \(E_D(u,v)\iff E_M(\mu(u),\mu(v))\), and
interpret class evaluation by the finite decoding just defined.  Induction
on equality-free formulas gives satisfaction transfer along \(\mu\),
using surjectivity at quantifiers.  Hence the represented \(R_b\) is
the pullback of \(R_b^M\).  For \(a\ne d\), choose \(W(s,a)\) to
falsify \(\varepsilon(P(z,s),a)\) when possible, and arbitrarily
otherwise.  At \(a=d\), every class instance is false, so the complemented
binder imposes no condition on \(W(s,d)\).  Use that free slice to choose
a member of \(\mu(s)\) whenever it has a member, lifting the selected
member by a leaf \(e_m\).  This is the inherited member-choice device;
\(V\) and \(W\) remain distinct operations.

The bad case makes every bad disjunct true.  In the hull case, EPSC closure
at empty inputs and hull members, the member-choice slice for nonempty
inputs, and the chosen support witnesses make all three hull clauses true.
Finally transport any countably infinite Boolean algebra onto the free
carrier with top equal to \(d\), interpret \(N\) as its NAND, and let
the two predicate-valued operations return its top or bottom according to
the just constructed relations.  All seven equations hold literally.
Since \(Q\) is a free tree constructor, \(Q(d)\ne d\), so the ground
query is separated.  This construction preserves the compatible class,
selector, and raw-membership roles of \(Q\); none is confused with the
Boolean truth value of the native tree itself.

\subsection{Relative calibration, evidence, and the exact charge}

Ground terms modulo the least congruence generated by all ground equation
instances form a countable model.  A ground equation is underivable
exactly when this model separates its endpoints.  The two constructions
above therefore prove, in the ordinary ZFC metatheory,
\[
 \mathcal E_{639}\nvdash Q(d)=d
 \quad\Longleftrightarrow\quad
 \mathcal E_{692}\nvdash N(d,d)=d.
\]
The right-hand side has the inspected report's inherited raw-EPSC/E713/ZFC
calibration.  Composing with it gives the intended equivalence with
\(\operatorname{Con}(\mathrm{ZFC})\).  The new relative theorem does not
require re-proving that older calibration, but its use is explicit.

The successive complete packets are
\begin{center}
\begin{tabular}{lrrr}
\hline
Packet & Equations & Equation nodes & Logical bits\\
\hline
Inspected E692 & 8 & 156 & 692\\
Self-descriptor E668 & 8 & 150 & 668\\
Truth deletion E643 & 7 & 144 & 643\\
Short query E639 & 7 & 144 & 639\\
\hline
\end{tabular}
\end{center}
The final accounting is exactly
\[
 \underbrace{43}_{\text{header}}
 +\underbrace{(7+4\cdot144+1)}_{\text{equations and terminator}}
 +\underbrace{12}_{\text{ground query}}
 =639.
\]

\texttt{build\_e639.py} checks the pinned E692 bytes, changes exactly the
three code projections in its last equation, and emits all three new
checkpoints with expanded equations and exact codec round trips.  It checks
the protected-tail and truth-recovery Boolean identities by symbolic
Boolean-ring normalization, and separately exercises the parameter-stack
coordinate calculation for lengths 0 through 64.  The latter is a finite
regression supporting the explicit head/tail induction; it is not an
exhaustive verification of the model construction.  The ordinary proof
above, the inherited Sh1 theorem, and the inherited raw-EPSC calibration
remain mathematical premises of the result.  No new Lean/Rocq endpoint,
internal PA proof-code translation, or binary-state count is supplied for
this alternative presentation.

The new contribution is the protected parameter-tail code, its compatible
model construction, and its exact composition with closure-based truth
recovery.  Source 104 retains credit for the earlier truth-equation
elimination on E946; Source 107 for the complemented binder and factoring;
the inspected E692 report for the structural-witness and selector
combination; and the published authors for Sh1.  No claim of historical
priority for general coding ideas, global bit minimality, or a binary
machine record follows from this packet.
